\documentclass[10pt,a4paper]{amsart}
\usepackage[margin=19mm]{geometry}
\usepackage{amsmath,amssymb,mathtools}
\usepackage[scr=rsfs,cal=euler]{mathalfa}
\usepackage{xfrac}
\usepackage[unicode,hidelinks,bookmarks=false]{hyperref}
\newcommand{\ie}{i.e.\ }
\newcommand{\cf}{cf.\ }
\newcommand{\resp}{respectively\ }
\newcommand{\Subsection}{\subsection}
\providecommand{\nobreakdash}{\nobreak\hbox{-}}
\setlength{\emergencystretch}{2em}
\multlinegap0pt
\usepackage{amssymb}
\usepackage{float}
\newtheorem{lem}{Lemma}
\newtheorem{thm}{Theorem}
\title{A homogeneous polynomial associated with cliques of graphs and its applications}
\author{Changjiang Bu, Jueru Liu, Haotian Zeng}
\date{Source: arXiv:2603.25365v3 (CC BY 4.0)}
\begin{document}
\maketitle

\noindent\emph{Source and licence.} A homogeneous polynomial associated with cliques of graphs and its applications. Version arXiv:2603.25365v3 (CC BY 4.0), distributed by arXiv under the Creative Commons Attribution 4.0 International licence, http://creativecommons.org/licenses/by/4.0/, as recorded in the arXiv licence metadata for this exact version and shown on the abstract page https://arxiv.org/abs/2603.25365v3. Full licence notice: \texttt{licenses/arxiv-2603.25365-CC-BY-4.0.txt}.\par\smallskip

\noindent\textit{Editorial scope.} Counted unit: the article's thm le (source0.tex lines 360--364). The article's complete original proof of it is delivered with it (source0.tex lines 366--400, one \texttt{proof} environment of 35 lines). No mathematics is written, repaired or re-derived: the statement and the proof are the article's own lines, in the article's own notation, and no retained line is substituted or re-flowed.  Retained uncounted as the source-local context the proof needs: lines 275--279.  Nothing was omitted from any retained span, so the package carries no omission record.  No retained line contains a citation, so the wrapper carries no bibliography and no bibliography entry.  No retained line contains a figure environment, an image-inclusion command or drawing code, so no image is delivered and the package declares no asset dependency.  Editorial formatting only, in addition: an AMS \texttt{amsart} preamble that restates the article's own declarations and macros used by the retained lines, namely the environments \texttt{lem}, \texttt{thm} on separate counters, together with the standard packages those lines need.  The retained environments print in this document as Lemma 1, Theorem 1 in order of appearance, whereas the article prints the same items with the numbering of its own full document.  The complete native source of the article as distributed in the arXiv e-print for this exact version is bundled unchanged as \texttt{sources/197204/source0.tex}.
\begin{lem}\label{w2}
Let $G$ be an $n$-vertex graph with clique number $\omega$ and let $2\le t\le\omega$. For any vector $x\in\mathbb{R}^{n}_{+}$ with $||x||_1=1$,
$$\sum_{I=\{i_1,i_2,\ldots,i_t\}\in C_t(G)} \frac{1}{t} \left(\sum_{j=1}^{t}\left(\alpha(i_j)\right)^t\binom{\alpha(i_j)}{t}^{-1}\right)x_I\le 1.$$
Equality holds if and only if the induced subgraph of $G$ on $\mathrm{supp}(x)$ is a complete $\omega$-partite graph with partition $V_1, V_2, \ldots, V_{\omega}$ satisfy $\sum_{v\in V_i}x_v=\frac{1}{\omega}$ for all $i\in[\omega]$.
\end{lem}

\begin{thm}\label{le}
Let $G$ be an $n$-vertex graph with clique number $\omega$ and let $2\le t\le\omega$. Then
$$\sum_{v\in V(G)}\frac{c_t(v)}{t}\sqrt[t-1]{\binom{\alpha(v)}{t}\left(\alpha(v)\right)^{-t}} \le \left( \sum_{v\in V(G)}\sqrt[t-1]{\binom{\alpha(v)}{t}\left(\alpha(v)\right)^{-t}} \right)^{t}.$$
Equality holds if and only if the graph obtained from $G$ by deleting edges not contained in $t$-cliques is a complete regular $\omega$-partite graph (possibly with some isolated vertices).
\end{thm}

\begin{proof}
A specific form of Muirhead's inequality states that for positive real number $z_1,z_1,\ldots,z_m$,
$$\sum_{k=1}^{m}\frac{1}{z_k}\le\left(\sum_{k=1}^{m}z_k^{m-1}\right)\frac{1}{\prod_{j=1}^{m}z_j},$$ and equality holds if and only if $z_1=\cdots=z_m$.
Then, we have
\begin{align*}
	 \sum_{v\in V(G)}\frac{c_t(v)}{t}\sqrt[t-1]{\frac{\binom{\alpha(v)}{t}}{\left(\alpha(v)\right)^t}}
= & \frac{1}{t} \sum_{\{i_1,i_2,\ldots,i_t\}\in C_t(G)}\left( \sum_{j=1}^{t} \sqrt[t-1]{\frac{\binom{\alpha(i_j)}{t}}{\left(\alpha(i_j)\right)^t}} \right)\\
\le & \frac{1}{t} \sum_{\{i_1,i_2,\ldots,i_t\}\in C_t(G)} \left( \left( \sum_{j=1}^{t} \frac{\left(\alpha(i_j)\right)^t}{\binom{\alpha(i_j)}{t}} \right) \prod_{j=1}^{t}\sqrt[t-1]{\frac{\binom{\alpha(i_j)}{t}}{\left(\alpha(i_j)\right)^t}} \right)\\
= & t \sum_{\{i_1,i_2,\ldots,i_t\}\in C_t(G)} \left( \left( \frac{1}{t^2} \sum_{j=1}^{t} \frac{\left(\alpha(i_j)\right)^t}{\binom{\alpha(i_j)}{t}} \right) \prod_{j=1}^{t}\sqrt[t-1]{\frac{\binom{\alpha(i_j)}{t}}{\left(\alpha(i_j)\right)^t}} \right).
\end{align*}
Let $V(G)=\{v_1,v_2,\ldots,v_n\}$. We construct a vector
$$x=\left( \sqrt[t-1]{\frac{\binom{\alpha(v_1)}{t}}{\left(\alpha(v_1)\right)^t}}, \sqrt[t-1]{\frac{\binom{\alpha(v_2)}{t}}{\left(\alpha(v_2)\right)^t}}, \ldots, \sqrt[t-1]{\frac{\binom{\alpha(v_n)}{t}}{\left(\alpha(v_n)\right)^t}} \right)^{\mathsf{T}}\in\mathbb{R}^{n}_{+}.$$
And let $\mathcal{W}(G)=(w_{i_1i_2\cdots i_t})$ be the weight $t$-clique tensor of $G$ with $w_{i_1i_2\cdots t_t}=\omega_{i_1i_2\cdots i_t}\frac{1}{(t-1)!}$ if $\{i_1,i_2,\ldots,i_t\}$ is a $t$-clique in $G$ and $w_{i_1i_2\cdots t_t}=0$ otherwise, where $\omega_{i_1i_2\cdots i_t}=\frac{1}{t^2} \sum_{j=1}^{t} \frac{\left(\alpha(i_j)\right)^t}{\binom{\alpha(i_j)}{t}}$.
Then
\begin{align*}
	\sum_{v\in V(G)}\frac{c_t(v)}{t}\sqrt[t-1]{\frac{\binom{\alpha(v)}{t}}{\left(\alpha(v)\right)^t}}\le \mathcal{W}(G)x^{t}.
\end{align*}
Let $y=\frac{x}{||x||_{1}}\in\mathbb{R}^{n}_{+}$. Then $||y||_{1}=\sum_{i=1}^{n}y_i=1$. By Lemma \ref{w2}, we have
$$\mathcal{W}(G)x^t=\mathcal{W}(G)y^t \cdot ||x||_{1}^{t}\le ||x||_{1}^{t}.$$
Thus,
$$	\sum_{v\in V(G)}\frac{c_t(v)}{t}\sqrt[t-1]{\frac{\binom{\alpha(v)}{t}}{\left(\alpha(v)\right)^t}} \le \mathcal{W}(G)x^{t} \le ||x||_{1}^{t} = \left( \sum_{v\in V(G)}\sqrt[t-1]{\frac{\binom{\alpha(v)}{t}}{\left(\alpha(v)\right)^t}} \right)^{t}.$$

Next, we consider the equality in the above inequality.
If the graph obtained from $G$ by deleting edges not contained in $t$-cliques is a complete regular $\omega$-partite graph, we will verify that equality holds.
Without loss of generality, assume that $\omega|n$ and $G$ is a complete regular $\omega$-partite graph.
Then $\alpha(v)=\omega$ and $c_t(v)=\binom{\omega-1}{t-1}\left( \frac{n}{\omega} \right)^{t-1}$ for any $v\in V(G)$.
Therefore, we can proof that the equality holds.

Conversely, if equality holds, then by Lemma \ref{w2}, $G[\mathrm{supp}(y)]$ is a complete $\omega$-partite graph with partition $V_1, V_2, \ldots, V_{\omega}$ satisfy $\sum_{v\in V_i}y_v=\frac{1}{\omega}$ for all $i\in[\omega]$.
By Muirhead's inequality, we know that $\alpha(v)=\omega$ for each $v$ contained in $t$-cliques.
Hence, for any $i\in[\omega]$,
$$\frac{1}{\omega}=\frac{1}{||x||_1}\sum_{v\in V_i}\sqrt[t-1]{\frac{\binom{\alpha(v)}{t}}{\left(\alpha(v)\right)^t}}=\frac{|V_i|}{||x||_1}\sqrt[t-1]{\frac{\binom{\omega}{t}}{\omega^t}},$$
i.e., $|V_1|=|V_2|=\cdots=|V_{\omega}|$.
Consequently, the graph obtained from $G$ by deleting edges not contained in $t$-cliques is a complete regular $\omega$-partite graph.
\end{proof}

\end{document}
